\subsection{Synthesis}

The measurement supports a filtering-dominant characterization of the
2026 blackout over a control-plane-withdrawal one: at every phase after
baseline, the dominant contributor to unreachability is routes that stay
announced while active probing collapses, not routes disappearing from BGP
(Section~\ref{sec:h1}). This holds at the level BGP visibility can actually
speak to across all eleven classification types (Section~\ref{sec:h3}), and
survives a control-population check that would catch a shared measurement
artifact (Sections~\ref{sec:h1}, \ref{sec:h3}). Within that broad picture,
restoration was selective by sector at the timescale the data can resolve,
though the ordering is not a clean single-axis priority list
(Section~\ref{sec:h3}): government and mobile restore fastest by median, but
mobile's tail is long, and the interesting selectivity signal -- the
state telecom operator sitting well ahead of the ordinary access-network
population -- is a coarser, three-tier structure (fast/government-and-
mobile, mid/state-infrastructure, slow/everything-else), not a fine-
grained ranking by sector.

\subsection{Disconfirming evidence}

Two results run against the cleanest version of this paper's own findings.
Both are kept in, not folded away. First, at P4 the control population's
own mean dark share (\controlDarkSharePFour{}) is higher than Iran's
(\darkSharePFour{}) -- restoration brings Iran's measured dark share back
down into the same noise floor the control population sits in throughout
the study, which is itself informative (Iran becomes statistically
indistinguishable from ``normal network noise'' by this specific metric
once restoration is underway) but means P4 cannot be read as ``still
clearly anomalous'' by the dark-share metric alone; the restoration-timing
and never-restored evidence in Section~\ref{sec:h3} is what actually
substantiates P4 as a distinct, partial-restoration regime. Second, H2's
originating claim -- that disconnection was executed at international-
gateway ASes specifically, as opposed to distributed across access networks
-- is not established by the aggregate upstream-transition-rate finding,
which is topology-blind; the one small, explicit check run against the two
named gateway ASNs was inconclusive on a sample far too small to draw a
conclusion from (Section~\ref{sec:h2}).

\subsection{Limitations}

\begin{itemize}
\item \textbf{H2's topological claim is untested at scale.} The aggregate
  churn-rate spike at the shutdown's onset is consistent with a structural
  reconfiguration but does not by itself distinguish gateway-level from
  distributed access-network disconnection. Breaking the upstream-change
  measurement down by each network's role in the topology -- gateway versus
  access network -- is the natural next analysis and was outside this
  study's scope.
\item \textbf{H4 has no control-population check.} The cross-window
  duration comparison (Section~\ref{sec:h4}) compares Iran to itself across
  five historical windows spanning 2019--2026, not to a control population
  at each of those times; shared measurement-infrastructure changes over
  that seven-year span (collector coverage, IODA methodology changes) are
  not ruled out the way H1/H3's within-study-period findings are.
\item \textbf{Classification coverage is partial.} \nClassified{} of
  \nIrAsns{} networks carry a category; the remainder are excluded from
  category-based analyses. Several categories have fewer than 20 members
  and carry count-only claims by design (Section~\ref{sec:methods}), since
  small groups are more sensitive to individual coding decisions than
  large ones.
\item \textbf{A small residual of control-population bins remain flagged.}
  \nControlArtifactBins{} of \nControlBins{} study-period bins cross the
  artifact threshold under the final calibrated dark-ratio threshold,
  concentrated late in the study period; this does not touch any Iranian
  finding in this paper (Section~\ref{sec:h1}) but was not chased to a root
  cause.
\end{itemize}

\subsection{Attribution discipline}

Every claim above about \emph{why} a pattern occurred, as opposed to
\emph{what} was measured, is bounded by the tiers in
Section~\ref{sec:methods}: this paper's measurement can establish that
routes stayed visible while reachability collapsed, and can time that
collapse relative to BGP events to the minute; it cannot, from BGP and
outage measurement alone, establish the internal decision process behind
that collapse. Claims in this paper about government action or intent are
sourced to external reporting, tiered no higher than \emph{consistent with
reporting}, and are never presented as conclusions the measurement itself
reaches.
